\section{Solution aqueuse d'acide butanoïque}
L'acide butanoïque de formule $\ce{C3H7CO2H}$, est l'un des composés
responsables de l'odeur très forte et du goût piquant de certains fromages et
beurres rances.  Il est présent dans les huiles végétales et les graisses
animales.

Cet exercice vise~:
\begin{itemize}
    \item l'étude d'une solution aqueuse d'acide butanoïque~;
    \item la détermination du pourcentage d'acide butanoïque dans un beurre.
\end{itemize}

\subsection{Étude de la solution aqueuse d'acide butanoïque}
On prépare, à $\qty{25}{\degreeCelsius}$, une solution aqueuse
$(S_{A})$ d'acide butanoïque de concentration molaire
$C_{A}=\qty{2,0e-3}{\M}$ et de volume $V_{A}=\qty{1,0}{\liter}$. La
mesure du $\mathrm{pH}$ de la solution $(S_{A})$ a donné la valeur
$\mathrm{pH}=3,76$.
\begin{enumerate}
    \item Écrire l'équation chimique modélisant la réaction de l'acide
          butanoïque avec l'eau.
    \item Dresser le tableau d'avancement de la réaction en utilisant
          les grandeurs $C_{A}$, $V_{A}$, l'avancement $x$ et
          l'avancement $x_{\eq}$ à l'état d'équilibre du système
          chimique.
    \item Déterminer la valeur de l'avancement maximal $x_{\max}$.
    \item Vérifier que la valeur de l'avancement à l'état d'équilibre
          est $x_{\eq}=\qty{1,74e-4}{\mole}$.
    \item Calculer la valeur du taux d'avancement final $\tau$. Que
          peut-on déduire~?
    \item Calculer la valeur de la constante d'équilibre $\mathrm{K}$
          associée à l'équation de cette réaction.
    \item Recopier sur votre copie le numéro de la question et écrire
          la lettre correspondante à la proposition vraie.

          Dans les conditions de l'expérience, la constante
          d'équilibre $\mathrm{K}$ associée à l'équation de la
          réaction~:
          \begin{minipage}{\linewidth}
              \begin{table}[H]
                  \centering
                  \setlength{\arrayrulewidth}{0.8pt}
                  \renewcommand{\arraystretch}{1.3}
                  \begin{tabularx}{\textwidth}{|
                        >{\arraybackslash\centering}p{7mm}|L|}
                      \hline 
                      A &  dépend à la fois de la composition initiale du
                            système chimique et de la température. \\
                      \hline 
                      B &  dépend uniquement de la composition initiale du
                            système chimique. \\
                      \hline 
                      C &  dépend uniquement du $\mathrm{pH}$ de la solution. \\
                      \hline 
                      D &  dépend uniquement de la température du système
                            chimique. \\
                      \hline 
                  \end{tabularx}
              \end{table}
          \end{minipage}
          
    \item Calculer la valeur du
          $\mathrm{pK_{A}}(\ce{C3H7CO2H_{(aq)}/C3H7CO$_2^-$_{(aq)}})$.
\end{enumerate}


          
\subsection{Détermination du pourcentage d'acide butanoïque dans un beurre}
Un beurre est rance si le pourcentage en masse d'acide butanoïque qu'il contient
est supérieur à $4\%$, c'est-à-dire qu'il y a plus de $\qty{4}{\gram}$ d'acide
butanoïque dans $\qty{100}{\gram}$ de beurre.

\begin{donnees}
    \begin{itemize}
        \item $M(\ce{C3H7CO2H})=\qty{88}{\gram\per\mol}$.
    \end{itemize}
\end{donnees}

Dans un bécher, on introduit $m_{b}=\qty{10,0}{\gram}$ de beurre fondu auquel on
ajoute de l'eau distillée. On agite afin de dissoudre dans l'eau la totalité de
l'acide butanoïque $\ce{C3H7CO2H}$ présent dans le beurre. On obtient une
solution aqueuse $(S)$ d'acide butanoïque de concentration molaire $C$ et de
volume $V_{0}=\qty{1,0}{\liter}$.

On dose le volume $V=\qty{10,0}{\mL}$ de la solution $(S)$ par une solution
aqueuse d'hydroxyde de sodium $\ce{Na^+_{(aq)} + HO^-_{(aq)}}$ de concentration
molaire $C_{B}=\qty{4,0e-3}{\M}$.
          
\begin{minipage}{0.60\linewidth}
    Le suivi pH-métrique du dosage permet d'obtenir la courbe
    ${\mathrm{pH}=f(V_{B})}$ (Figure ci-contre).  On considère que seul l'acide
    butanoïque réagit avec le réactif titrant.
    \begin{enumerate}
        \item Écrire l'équation de la réaction du dosage sachant
              qu'elle est totale.
        \item Déterminer graphiquement la valeur du volume $V_{B,\mathrm{E}}$ à
              l'équivalence.
        \item Calculer la valeur de $C$.
        \item Déterminer la masse de l'acide butanoïque présent dans
              la masse $m_{b}=\qty{10,0}{\gram}$ du beurre. Le beurre
              étudié est-il rance~? Justifier.
    \end{enumerate}
\end{minipage}
\hfill
\begin{minipage}{0.38\linewidth}
    \flushright
    \begin{figure}[H]
        \centering
        \includegraphics[width=\textwidth]{2026-03-23_122336-}
    \end{figure}
\end{minipage}



